\documentclass[10pt]{article}
\usepackage{style}
\renewcommand{\examround}{1학기 기말 예상 — 2회 답지}
\renewcommand{\examvol}{vol.2}

\setlength{\columnsep}{8mm}
\setlength{\columnseprule}{0pt}

\usepackage{eso-pic}
\AddToShipoutPictureFG{%
  \ifnum\value{page}=1
    \begin{tikzpicture}[overlay,remember picture]
      \draw[color=vB-rule,line width=0.4pt]
        ([xshift=1mm,yshift=-125.6mm]current page.north)
        -- ([xshift=1mm,yshift=15.8mm]current page.south);
    \end{tikzpicture}%
  \else
    \begin{tikzpicture}[overlay,remember picture]
      \draw[color=vB-rule,line width=0.4pt]
        ([xshift=1mm,yshift=-13.6mm]current page.north)
        -- ([xshift=1mm,yshift=15.8mm]current page.south);
    \end{tikzpicture}%
  \fi
}

\begin{document}

\answerheader{1학기 기말 예상 — 2회 풀이}{공통수학1 \textperiodcentered\ 다항식 \textperiodcentered\ 여러 가지 방정식과 부등식 \textperiodcentered\ 순열과 조합 \textperiodcentered\ 행렬}

\begin{quickgrid}
\quick{1}{⑤ $60$} & \quick{2}{③ $30$} & \quick{3}{② $\tfrac{2}{3}$} & \quick{4}{① $4$} & \quick{5}{② $11$} \\
\quick{6}{④ $13$} & \quick{7}{③ $3$} & \quick{8}{④ $7$} & \quick{9}{⑤ $14400$} & \quick{10}{⑤ $4$} \\
\quick{11}{④ $6$} & \quick{12}{③ $5$} & \quick{13}{③ $12$} & \quick{14}{② $k>1$} & \quick{15}{⑤ $27$} \\
\quick{16}{⑤ $200$} & \quick{17}{③ $336$} & \quick{18}{④ $576$} & \quick{19}{④ $10$} & \quick{20}{③ $\tfrac{4}{5}DC$의 $(1,2)$ 성분} \\
\quick{21}{③ $4$} & \quick{22}{③ $-3$} & \quick{23}{① $0$} & & \\
\end{quickgrid}

\examsection{객관식 풀이}

\raggedcolumns
\begin{multicols}{2}

% ─────────────── #1 ───────────────
\soltitle{1}{대칭성으로 반씩 다루기}{2}
\begin{answerbox}정답: ⑤ $60$\end{answerbox}

$5$명을 일렬로 세우는 전체 방법의 수는
\[ 5! = 120 \]
이다. 임의의 두 명 A, B를 고정하면 A가 B보다 앞에 서는 경우와 뒤에 서는 경우는 대칭이므로 각각 전체의 절반을 차지한다. 그러므로 구하는 방법의 수는
\[ \tfrac{1}{2} \times 120 = 60. \]

% ─────────────── #2 ───────────────
\soltitle{2}{전체에서 조건 어긋난 경우 빼기}{2}
\begin{answerbox}정답: ③ $30$\end{answerbox}

전체 $7$명 중 대표 $3$명을 뽑는 방법의 수는
\[ {}_7\mathrm{C}_3 = 35. \]

$3$명 모두 남학생인 경우는 ${}_4\mathrm{C}_3 = 4$가지, $3$명 모두 여학생인 경우는 ${}_3\mathrm{C}_3 = 1$가지이다.

남학생과 여학생이 모두 포함되는 경우의 수는 전체에서 남학생만 뽑는 경우와 여학생만 뽑는 경우를 빼서
\[ 35 - 4 - 1 = 30. \]

% ─────────────── #3 ───────────────
\soltitle{3}{절댓값을 포함한 방정식 다루기}{2}
\begin{answerbox}정답: ② $\tfrac{2}{3}$\end{answerbox}

$|x - 3| = 2x + 1$의 우변이 음이 아닐 조건 $2x + 1 \ge 0$, 즉 $x \ge -\tfrac{1}{2}$이 성립해야 한다.

$(\!$ⅰ$)$ $x \ge 3$일 때 $x - 3 = 2x + 1$에서 $x = -4$. 조건 $x \ge 3$에 모순.

$(\!$ⅱ$)$ $-\tfrac{1}{2} \le x < 3$일 때 $-(x - 3) = 2x + 1$에서 $3 - x = 2x + 1$, $3x = 2$, $x = \tfrac{2}{3}$. 조건을 만족한다.

그러므로 모든 실수 $x$의 값의 합은
\[ \tfrac{2}{3}. \]

% ─────────────── #4 ───────────────
\soltitle{4}{행렬의 곱 계산 다루기}{2}
\begin{answerbox}정답: ① $4$\end{answerbox}

$A = \begin{pmatrix} 1 & 2 \\ 3 & -1 \end{pmatrix}$이므로
\[ A^2 = \begin{pmatrix} 1 + 6 & 2 - 2 \\ 3 - 3 & 6 + 1 \end{pmatrix} = \begin{pmatrix} 7 & 0 \\ 0 & 7 \end{pmatrix}, \]
\[ 2A = \begin{pmatrix} 2 & 4 \\ 6 & -2 \end{pmatrix}. \]
그러므로
\[ A^2 - 2A = \begin{pmatrix} 5 & -4 \\ -6 & 9 \end{pmatrix}, \]
모든 성분의 합은
\[ 5 + (-4) + (-6) + 9 = 4. \]

% ─────────────── #5 ───────────────
\soltitle{5}{조합의 성질 다루기}{2}
\begin{answerbox}정답: ② $11$\end{answerbox}

${}_n\mathrm{C}_r = \dfrac{n!}{r! \, (n-r)!}$이므로
\[ 2 \cdot {}_n\mathrm{C}_3 = {}_n\mathrm{C}_4 \iff 2 \cdot \dfrac{n!}{3! \, (n-3)!} = \dfrac{n!}{4! \, (n-4)!}. \]
양변을 $\dfrac{n!}{4! \, (n-4)!}$로 나누면
\[ 2 \cdot \dfrac{4!}{3! \, (n-3)} = 1, \quad \dfrac{8}{n - 3} = 1. \]
그러므로 $n - 3 = 8$, 즉
\[ n = 11. \]

% ─────────────── #6 ───────────────
\soltitle{6}{나머지정리로 두 조건 다루기}{2}
\begin{answerbox}정답: ④ $13$\end{answerbox}

$P(x) = x^3 + ax + b$라 하자. 나머지정리에 의해
\begin{align*}
P(1) = 1 + a + b &= 2, \\
P(-1) = -1 - a + b &= -6.
\end{align*}
즉
\[ a + b = 1, \quad -a + b = -5. \]
두 식을 더하면 $2b = -4$, $b = -2$. 첫째 식에 대입하여 $a = 3$이다. 그러므로
\[ a^2 + b^2 = 9 + 4 = 13. \]

% ─────────────── #7 ───────────────
\soltitle{7}{대칭식으로 값 구하기}{2}
\begin{answerbox}정답: ③ $3$\end{answerbox}

$x^3 + y^3 = (x + y)^3 - 3xy(x + y)$이므로
\[ 28 = 4^3 - 3xy \cdot 4 = 64 - 12\, xy. \]
정리하면
\[ 12\, xy = 36, \quad xy = 3. \]

% ─────────────── #8 ───────────────
\soltitle{8}{절댓값 부등식 정수해 다루기}{3}
\begin{answerbox}정답: ④ $7$\end{answerbox}

절댓값 안의 부호에 따라 세 구간으로 나누어 푼다.

$(\!$ⅰ$)$ $x \ge \tfrac{3}{2}$일 때 부등식은
\[ (2x - 3) + (x + 1) = 3x - 2 \le 10, \]
즉 $x \le 4$이므로 $\tfrac{3}{2} \le x \le 4$. 정수 $x$는 $2, 3, 4$의 $3$개.

$(\!$ⅱ$)$ $-1 \le x < \tfrac{3}{2}$일 때 부등식은
\[ (3 - 2x) + (x + 1) = 4 - x \le 10, \]
즉 $x \ge -6$이므로 $-1 \le x < \tfrac{3}{2}$ 전체. 정수 $x$는 $-1, 0, 1$의 $3$개.

$(\!$ⅲ$)$ $x < -1$일 때 부등식은
\[ (3 - 2x) - (x + 1) = 2 - 3x \le 10, \]
즉 $x \ge -\tfrac{8}{3}$이므로 $-\tfrac{8}{3} \le x < -1$. 정수 $x$는 $-2$의 $1$개.

이상에서 조건을 만족시키는 정수 $x$의 개수는
\[ 3 + 3 + 1 = 7. \]

% ─────────────── #9 ───────────────
\soltitle{9}{이웃하지 않는 순열 다루기}{2}
\begin{answerbox}정답: ⑤ $14400$\end{answerbox}

먼저 남학생 $5$명을 일렬로 배열한다. 배열의 수는
\[ 5! = 120. \]
그 뒤 남학생 사이와 양 끝의 $6$자리 중 서로 다른 $3$자리를 골라 여학생 $3$명을 배치하면 여학생끼리 서로 이웃하지 않게 된다. 이 방법의 수는
\[ {}_6\mathrm{P}_3 = 6 \times 5 \times 4 = 120. \]
곱의 법칙에 의하여 구하는 방법의 수는
\[ 120 \times 120 = 14400. \]

% ─────────────── #10 ───────────────
\soltitle{10}{이차부등식 정수해 조건 다루기}{3}
\begin{answerbox}정답: ⑤ $4$\end{answerbox}

$x^2 - (a + 2)x + 2a = (x - 2)(x - a)$이므로 부등식은
\[ (x - 2)(x - a) < 0. \]
$2$와 $a$의 크기 관계에 따라 두 경우로 나눈다.

$(\!$ⅰ$)$ $a > 2$일 때 해는 $2 < x < a$. 정수 $x$는 $3, 4, 5, \ldots$이고 개수가 $3$이 되려면 $3, 4, 5$가 모두 해에 속하고 $6$은 속하지 않아야 하므로
\[ 5 < a \le 6. \]

$(\!$ⅱ$)$ $a < 2$일 때 해는 $a < x < 2$. 정수 $x$는 $\ldots, -1, 0, 1$이고 개수가 $3$이 되려면 $-1, 0, 1$이 모두 해에 속하고 $-2$는 속하지 않아야 하므로
\[ -2 \le a < -1. \]

$(\!$ⅲ$)$ $a = 2$일 때 $(x - 2)^2 < 0$은 해가 존재하지 않는다.

두 경우를 합하면 실수 $a$의 최댓값은 $6$, 최솟값은 $-2$이므로 그 합은
\[ 6 + (-2) = 4. \]

% ─────────────── #11 ───────────────
\soltitle{11}{삼차방정식 인수분해로 실근 조건 다루기}{3}
\begin{answerbox}정답: ④ $6$\end{answerbox}

$x^3 + (8 - a)x^2 + (a^2 - 8a)x - a^3$에 $x = a$를 대입하면
\begin{align*}
&a^3 + (8-a)a^2 + (a^2 - 8a)a - a^3 \\
&\quad = a^3 + 8a^2 - a^3 + a^3 - 8a^2 - a^3 = 0
\end{align*}
이므로 $x = a$는 항상 이 방정식의 근이다. 조립제법 또는 다항식의 나눗셈으로 인수분해하면
\begin{align*}
&x^3 + (8-a)x^2 + (a^2-8a)x - a^3 \\
&\quad = (x - a)(x^2 + 8x + a^2).
\end{align*}

서로 다른 세 실근을 가지려면 이차식 $x^2 + 8x + a^2 = 0$이 서로 다른 두 실근을 가지며, 그 두 근 모두 $a$와 달라야 한다.

$(\!$ⅰ$)$ 판별식 조건:
\[ D/4 = 16 - a^2 > 0 \;\Rightarrow\; -4 < a < 4. \]

$(\!$ⅱ$)$ $x = a$가 이차식의 근이 아니어야 하므로
\[ a^2 + 8a + a^2 = 2a(a+4) \ne 0, \]
즉 $a \ne 0$이고 $a \ne -4$이다.

$-4 < a < 4$에서 $a \ne 0$인 정수 $a$는
\[ -3,\; -2,\; -1,\; 1,\; 2,\; 3 \]
의 $6$개이다.

% ─────────────── #12 ───────────────
\soltitle{12}{판별식으로 항상 성립 조건 다루기}{2}
\begin{answerbox}정답: ③ $5$\end{answerbox}

모든 실수 $x$에 대하여 $x^2 + (a - 1)x + a > 0$이 성립할 필요충분조건은 좌변의 판별식이 음수인 것이다.
\[ D = (a - 1)^2 - 4a = a^2 - 6a + 1 < 0. \]
$a^2 - 6a + 1 = 0$의 해는 $a = 3 \pm 2\sqrt{2}$이므로
\[ 3 - 2\sqrt{2} < a < 3 + 2\sqrt{2}. \]
$2\sqrt{2} \approx 2.83$이므로 대략 $0.17 < a < 5.83$이 되어 정수 $a$는
\[ 1,\; 2,\; 3,\; 4,\; 5 \]
의 $5$개이다.

% ─────────────── #13 ───────────────
\soltitle{13}{행렬의 곱셈 조건 다루기}{2}
\begin{answerbox}정답: ③ $12$\end{answerbox}

$A = \begin{pmatrix} a & 1 \\ 2 & 3 \end{pmatrix}$, $B = \begin{pmatrix} 1 & 2 \\ b & 1 \end{pmatrix}$의 곱을 각각 계산하면
\begin{align*}
AB &= \begin{pmatrix} a + b & 2a + 1 \\ 2 + 3b & 7 \end{pmatrix}, \\
BA &= \begin{pmatrix} a + 4 & 7 \\ ab + 2 & b + 3 \end{pmatrix}.
\end{align*}

$AB = BA$가 성립하려면 성분끼리 일치해야 한다. $(1, 1)$ 성분에서
\[ a + b = a + 4 \;\Rightarrow\; b = 4. \]
$(2, 2)$ 성분에서
\[ 7 = b + 3 \;\Rightarrow\; b = 4 \]
로 확인된다. $(2, 1)$ 성분에서
\[ 2 + 3b = ab + 2 \;\Rightarrow\; 3b = ab \;\Rightarrow\; b(a - 3) = 0. \]
$b = 4 \ne 0$이므로 $a = 3$이다.

따라서
\[ ab = 3 \times 4 = 12. \]

% ─────────────── #14 ───────────────
\soltitle{14}{사차방정식을 이차방정식으로 치환하기}{2}
\begin{answerbox}정답: ④ $k > 9$\end{answerbox}

$t = x^2$로 놓으면 $x$가 실수일 조건은 $t \ge 0$이고, 주어진 방정식은 다음과 같이 $t$에 대한 이차방정식으로 바뀐다.
\[ f(t) = t^2 - 6t + k = 0 \quad(t \ge 0) \]
$f(t) = 0$의 두 근을 $t_1$, $t_2$라 하면 근과 계수의 관계로
\[ t_1 + t_2 = 6,\quad t_1 t_2 = k \]
이다. 두 근의 합이 양수이므로 두 근이 모두 음수가 되는 경우는 없다. 따라서 $x$에 대한 사차방정식이 실근을 갖지 않는 조건은 $f(t) = 0$이 실근 자체를 갖지 않는 것과 같다. 판별식
\[ 36 - 4k < 0 \;\Rightarrow\; k > 9. \]

% ─────────────── #15 ───────────────
\soltitle{15}{삼차방정식 인수분해로 세 근 다루기}{2}
\begin{answerbox}정답: ⑤ $27$\end{answerbox}

좌변을 인수분해한다.
\begin{align*}
&x^3 - 3x^2 - 4x + 12 \\
&\quad = x^2(x - 3) - 4(x - 3) \\
&\quad = (x - 3)(x^2 - 4) \\
&\quad = (x - 3)(x - 2)(x + 2).
\end{align*}

그러므로 세 근은 $3$, $2$, $-2$이다. 따라서
\[ \alpha^3 + \beta^3 + \gamma^3 = 3^3 + 2^3 + (-2)^3 = 27 + 8 - 8 = 27. \]

% ─────────────── #16 ───────────────
\soltitle{16}{나머지정리 확장 — 세제곱 관계}{3}
\begin{answerbox}정답: ⑤ $200$\end{answerbox}

$x^2 + x + 1 = 0$의 한 허근을 $\omega$라 하자. $(x - 1)(x^2 + x + 1) = x^3 - 1$이므로
\[ \omega^3 - 1 = 0,\quad \text{즉}\ \omega^3 = 1. \]

$x^{50} + 100x + 3$을 $x^2 + x + 1$로 나눈 몫을 $Q(x)$, 나머지를 $R(x) = ax + b$라 하면
\[ x^{50} + 100x + 3 = (x^2 + x + 1)\,Q(x) + ax + b. \]

이 등식에 $x = \omega$를 대입하면 $\omega^2 + \omega + 1 = 0$이므로 $\omega^{50} + 100\omega + 3 = a\omega + b$. $\omega^{50} = (\omega^3)^{16} \cdot \omega^2 = \omega^2$이고 $\omega^2 = -1 - \omega$이므로
\[ \omega^{50} + 100\omega + 3 = (-1 - \omega) + 100\omega + 3 = 2 + 99\omega. \]

따라서 $a\omega + b = 99\omega + 2$이고 $\omega$가 허수이므로 $a = 99$, $b = 2$이다. $R(x) = 99x + 2$이므로
\[ R(2) = 198 + 2 = 200. \]

% ─────────────── #17 ───────────────
\soltitle{17}{전체에서 이웃한 경우 빼기·포함배제}{3}
\begin{answerbox}정답: ③ $336$\end{answerbox}

$6$장의 카드를 일렬로 나열하는 전체 경우의 수는
\[ 6! = 720. \]

조건에 어긋나는 두 가지 경우를 각각 살핀다.

$(\!$ⅰ$)$ $1$과 $2$의 카드가 서로 이웃하는 경우: $1, 2$를 하나의 묶음으로 취급하면 그 안에서 두 순서 $2$가지, 나머지 $5$개 대상 배열은 $5!$가지이므로
\[ 2 \times 5! = 240. \]

$(\!$ⅱ$)$ $3$과 $4$의 카드가 서로 이웃하는 경우: 마찬가지로 $240$가지.

$(\!$ⅲ$)$ 두 조건이 동시에 성립하는 경우: $1, 2$와 $3, 4$를 각각 묶음으로 취급하면 두 묶음의 내부 순서 $2 \times 2$가지, 나머지 $4$개 대상 배열은 $4!$가지이므로
\[ 2 \times 2 \times 4! = 96. \]

두 조건 중 적어도 하나가 성립하는 경우의 수는
\[ 240 + 240 - 96 = 384. \]

따라서 두 조건 모두 어긋나지 않는 경우의 수는 전체에서 위 값을 빼서
\[ 720 - 384 = 336. \]

% ─────────────── #18 ───────────────
\soltitle{18}{짝수·홀수 카드 번갈아 배치 다루기}{3}
\begin{answerbox}정답: ④ $576$\end{answerbox}

두 이웃한 카드에 적힌 수의 합이 홀수라는 조건은 짝수와 홀수 카드가 서로 번갈아 놓인다는 조건과 같다. 카드 $1$$\sim$$8$ 중 홀수는 $1, 3, 5, 7$의 $4$개, 짝수는 $2, 4, 6, 8$의 $4$개이다.

$5$장을 뽑아 나열하므로 배치는 짝-홀-짝-홀-짝 또는 홀-짝-홀-짝-홀의 두 가지 형태이다.

$(\!$ⅰ$)$ 짝-홀-짝-홀-짝 형태: 짝수 $3$장을 순서대로 배치하는 방법은 ${}_4\mathrm{P}_3 = 24$가지, 홀수 $2$장을 순서대로 배치하는 방법은 ${}_4\mathrm{P}_2 = 12$가지이므로
\[ 24 \times 12 = 288. \]

$(\!$ⅱ$)$ 홀-짝-홀-짝-홀 형태: 마찬가지로 홀수 $3$장 ${}_4\mathrm{P}_3 = 24$가지, 짝수 $2$장 ${}_4\mathrm{P}_2 = 12$가지이므로
\[ 24 \times 12 = 288. \]

합의 법칙에 의하여
\[ 288 + 288 = 576. \]

% ─────────────── #19 ───────────────
\soltitle{19}{$1$부터 $k$까지 자연수의 합 이용}{2}
\begin{answerbox}정답: ④ $10$\end{answerbox}

$A + A^2 + \cdots + A^k$의 $(1, 2)$ 성분은 각 $A^n$의 $(1, 2)$ 성분 $n$을 모두 더한 값이다. 그러므로
\[ S = 1 + 2 + 3 + \cdots + k = \tfrac{k(k+1)}{2}. \]

$S = 55$이므로
\[ k(k+1) = 110. \]
연속한 두 자연수의 곱이 $110$인 경우는 $10 \times 11 = 110$뿐이므로
\[ k = 10. \]

% ─────────────── #20 ───────────────
\soltitle{20}{행렬 곱의 성분 해석하기}{3}
\begin{answerbox}정답: ③ $\tfrac{4}{5}DC$의 $(1, 2)$ 성분\end{answerbox}

A 형 축구공의 판매량은
\[ 500 \times 0.8 = 400 \; (\text{만 개}), \]
B 형 축구공의 판매량은 $300 \times 0.8 = 240$(만 개)이다. 판매가는 A 형이 $5$, B 형이 $7$이므로 이번 달 총 매출액은
\begin{align*}
&400 \times 5 + 240 \times 7 \\
&\quad = \tfrac{4}{5}(500 \times 5 + 300 \times 7)\quad (\text{만 원}\cdot\text{만 개}).
\end{align*}

한편 $DC$의 $(1, 2)$ 성분은 $D$의 $1$행(생산량 행 $(500,\,300)$)과 $C$의 $2$열(판매가 열 $\begin{pmatrix} 5 \\ 7 \end{pmatrix}$)의 곱으로
\[ (DC)_{12} = 500 \cdot 5 + 300 \cdot 7. \]

따라서 총 매출액은 $\dfrac{4}{5} DC$의 $(1, 2)$ 성분이다.

% ─────────────── #21 ───────────────
\soltitle{21}{이차함수와 절댓값 그래프의 교점 개수 다루기}{3}
\begin{answerbox}정답: ③ $4$\end{answerbox}

$g(x) = |x - 2|$의 정의에 따라 $x \ge 2$와 $x < 2$의 두 경우로 나누어
\[ f(x) = g(x) \]
의 실근을 센다.

$(\!$ⅰ$)$ $x \ge 2$: $g(x) = x - 2$이므로
\[ x^2 - 4x + a = x - 2 \;\Longleftrightarrow\; x^2 - 5x + (a + 2) = 0. \tag{1} \]

$(\!$ⅱ$)$ $x < 2$: $g(x) = 2 - x$이므로
\[ x^2 - 4x + a = 2 - x \;\Longleftrightarrow\; x^2 - 3x + (a - 2) = 0. \tag{2} \]

두 이차식의 판별식은 모두 $17 - 4a$로 같다. 판별식 $\ge 0$의 조건은 $a \le \dfrac{17}{4}$.

(1)의 두 근을 $x_1 \le x_2$라 하면 $x_1 + x_2 = 5$, $x_1 x_2 = a + 2$이고 두 근 모두 $x \ge 2$인 조건은
\[ x_1 \ge 2 \;\Longleftrightarrow\; \sqrt{17 - 4a} \le 1 \;\Longleftrightarrow\; a \ge 4. \]

(2)의 두 근을 $y_1 \le y_2$라 하면 $y_1 + y_2 = 3$이고 두 근 모두 $x < 2$인 조건은
\[ y_2 < 2 \;\Longleftrightarrow\; \sqrt{17 - 4a} < 1 \;\Longleftrightarrow\; a > 4. \]

따라서 유효 교점의 개수는 다음과 같다.
\begin{itemize}
\item $a < 4$: (1)에서 $x_2$만 유효, (2)에서 $y_1$만 유효 → $2$개.
\item $a = 4$: (1)의 근 $x = 2,\, 3$ 모두 유효, (2)의 근 $x = 1,\, 2$ 중 $x = 1$만 유효. 서로 다른 교점은 $x = 1,\, 2,\, 3$의 $3$개.
\item $4 < a < \tfrac{17}{4}$: (1)·(2) 각각 두 근 모두 유효 → $4$개.
\item $a = \tfrac{17}{4}$: 중근 각 $1$개씩 → $2$개.
\item $a > \tfrac{17}{4}$: 실근 없음.
\end{itemize}

서로 다른 세 점에서 만나는 경우는 $a = 4$뿐이다.

% ─────────────── #22 ───────────────
\soltitle{22}{이차함수 조건 통합과 근의 위치 판정}{3}
\begin{answerbox}정답: ③ $-3$\end{answerbox}

조건 $(\!$다$)$에서 $f(x) \le 0$의 해가 $\alpha \le x \le \beta$ 꼴이므로 $f(x) = 0$은 서로 다른 두 실근 $\alpha$, $\beta$($\alpha < \beta$)를 가지며
\[ f(x) = (x - \alpha)(x - \beta),\qquad \beta - \alpha = 4. \]

조건 $(\!$가$)$ $f(-1) = 0$에서 $-1$이 두 근 중 하나이므로 다음 두 가지 경우가 있다.
\begin{itemize}
\item 경우 A: $\alpha = -1$, $\beta = 3$. 즉 $f(x) = (x + 1)(x - 3)$.
\item 경우 B: $\alpha = -5$, $\beta = -1$. 즉 $f(x) = (x + 5)(x + 1)$.
\end{itemize}

각 경우에 대해 조건 $(\!$나$)$의 방정식 $f(x) = 3$의 두 실근이 부등식 $-2 < x < 5$를 모두 만족시키는지 확인한다.

경우 A: $f(x) = x^2 - 2x - 3$. $f(x) = 3$에서
\[ x^2 - 2x - 6 = 0 \;\Longrightarrow\; x = 1 \pm \sqrt{7}. \]
$\sqrt{7}$은 $2$와 $3$ 사이의 수이므로 두 근은 $1 - \sqrt{7} \approx -1.65$, $1 + \sqrt{7} \approx 3.65$이고 둘 다 $-2 < x < 5$에 속한다. 조건 $(\!$나$)$ 만족.

경우 B: $f(x) = x^2 + 6x + 5$. $f(x) = 3$에서
\[ x^2 + 6x + 2 = 0 \;\Longrightarrow\; x = -3 \pm \sqrt{7}. \]
두 근 중 $-3 - \sqrt{7} \approx -5.65$는 $-2 < x < 5$에 속하지 않는다. 조건 $(\!$나$)$ 위반.

따라서 경우 A만 유효이고 $f(x) = (x + 1)(x - 3)$이다. 그러므로
\[ f(2) = (2 + 1)(2 - 3) = 3 \cdot (-1) = -3. \]

% ─────────────── #23 ───────────────
\soltitle{23}{절댓값이 있는 방정식의 실근 개수 다루기}{3}
\begin{answerbox}정답: ① $0$\end{answerbox}

$f(x) = 0$에서 $|x^2 - tx| = 4$이므로
\[ x^2 - tx = 4 \quad \text{또는} \quad x^2 - tx = -4. \]

두 이차방정식의 판별식은 각각
\begin{align*}
D_1 &= t^2 + 16, \\
D_2 &= t^2 - 16.
\end{align*}

$D_1$은 항상 양수이므로 $x^2 - tx - 4 = 0$은 항상 서로 다른 두 실근을 갖는다. 또 두 식 $x^2 - tx = 4$와 $x^2 - tx = -4$가 공통근을 가지면 $4 = -4$가 되어 모순이므로 두 식의 실근은 서로 다르다.

$(\!$ⅰ$)$ $|t| < 4$: $D_2 < 0$이므로 $x^2 - tx + 4 = 0$은 실근을 가지지 않는다. 서로 다른 실근의 총 개수는 $2$.

$(\!$ⅱ$)$ $|t| = 4$: $D_2 = 0$이므로 $x^2 - tx + 4 = 0$은 중근 $x = \tfrac{t}{2}$를 가진다. 서로 다른 실근의 총 개수는 $2 + 1 = 3$.

$(\!$ⅲ$)$ $|t| > 4$: $D_2 > 0$이므로 서로 다른 두 실근을 갖는다. 서로 다른 실근의 총 개수는 $2 + 2 = 4$.

$g(t) \ge 3$이 되는 조건은 $|t| \ge 4$, 즉
\[ t \le -4 \quad \text{또는} \quad t \ge 4. \]

따라서 $\alpha = -4$, $\beta = 4$이므로
\[ \alpha + \beta = 0. \]

\end{multicols}

\end{document}
